\subsection{有限族代数的张量积}

A 恒表示一个含单位元的交换环。设 \((\mathbf{E}_i)_{i\in I}\) 是一个有限的 A-代数族，令 \(\mathbf{E} = \bigotimes_{i\in I}\mathbf{E}_i\) 是诸 A-模 \(\mathbf{E}_i\) 的张量积 A-模（II, § 3, 第 9 号）。我们要在 E 上定义一个 A-代数结构。设 \(m_i\colon \mathbf{E}_i\otimes_\mathbf{A}\mathbf{E}_i\to \mathbf{E}_i\) 是定义 \(\mathbf{E}_i\) 上乘法的 A-线性映射（§1, 第 3 号）。考虑 A-线性映射

\[
m ^ {\prime} = \bigotimes_ {i \in I} m _ {i}: \bigotimes_ {i \in I} (E _ {i} \otimes_ {A} E _ {i}) \rightarrow \bigotimes_ {i \in I} E _ {i} = E;
\]

复合映射

\[
\left(\bigotimes_ {i \in I} \mathrm{E} _ {i}\right) \otimes_ {\mathrm{A}} \left(\bigotimes_ {i \in I} \mathrm{E} _ {i}\right) \xrightarrow {\tau} \bigotimes_ {i \in I} \left(\mathrm{E} _ {i} \otimes_ {\mathrm{A}} \mathrm{E} _ {i}\right) \xrightarrow {m ^ {\prime}} \bigotimes_ {i \in I} \mathrm{E} _ {i}
\]

（其中 \(\tau\) 是结合性同构（II, §3, 第 9 号））是一个 A-线性映射 \(m\colon \mathbf{E} \otimes_{\mathbf{A}} \mathbf{E} \to \mathbf{E}\) ；我们将看到 \(m\) 在 \(\mathbf{E}\) 上定义了一个（结合且含单位元的）代数结构。事实上，具体算出由 \(m\) 定义的乘法，就得到公式

\[
\left(\bigotimes_ {i \in I} x _ {i}\right) \left(\bigotimes_ {i \in I} y _ {i}\right) = \bigotimes_ {i \in I} \left(x _ {i} y _ {i}\right) \quad \text { for } x _ {i}, y _ {i} \text { in } E _ {i} \text { and } i \in I.\tag{1}
\]

于是已经可以由线性看出，若 \(e_i\) 是 \(\mathbf{E}_i\) 的单位元，则 \(e = \bigotimes_{i \in I} e_i\) 是 \(\mathbf{E}\) 的单位元。另一方面，每个 \(\mathbf{E}_i\) 的结合性蕴含关系

\[
\left(\left(\bigotimes_ {i \in I} x _ {i}\right) \left(\bigotimes_ {i \in I} y _ {i}\right)\right) \left(\bigotimes_ {i \in I} z _ {i}\right) = \bigotimes_ {i \in I} \left(x _ {i} y _ {i} z _ {i}\right) = \left(\bigotimes_ {i \in I} x _ {i}\right) \left(\left(\bigotimes_ {i \in I} y _ {i}\right) \left(\bigotimes_ {i \in I} z _ {i}\right)\right)
\]

从而由线性性质得到关系 \(x(yz) = (xy)z\) 对 E 中所有 x, y, z 成立。

定义 1. 给定 A 上代数的一个族 \((\mathbf{E}_i)_{i\in I}\) ，该族的张量积（记作 \(\bigotimes_{i\in I}\mathbf{E}_i\) ，或者，当 I 是区间 \([1,n]\) （N 中的）时，记作 \(\mathbf{E}_1\otimes_\mathbf{A}\mathbf{E}_2\otimes \dots \otimes_\mathbf{A}\mathbf{E}_n\) ，或简记作 \(\mathbf{E}_1\otimes \mathbf{E}_2\otimes \dots \otimes \mathbf{E}_n\) ）是赋予诸 A-模 \(\mathbf{E}_i\) 的张量积以由 (1) 定义的乘法而得到的代数。

关系式 (1) 表明 \(\bigotimes_{i\in I}E_i^0\) （诸 \(E_{i}\) 的反代数的张量积）是 \(\bigotimes_{i\in I}E_{i}\) 的反代数；特别地，若诸 \(E_{i}\) 都是交换的，则 \(\bigotimes_{i\in I}E_{i}\) 也是交换的 .

设 \((\mathbf{E}_i)_{i\in I}\) 和 \((\mathbf{F}_i)_{i\in I}\) 是具有同一有限指标集 I 的两个 A-代数族。对每个 \(i\in I\) ，设 \(f_{i}\colon \mathbf{E}_{i}\to \mathbf{F}_{i}\) 是一个 A-代数同态。则 A-线性映射

\[
f = \bigotimes_ {i \in I} f _ {i}: \bigotimes_ {i \in I} E _ {i} \rightarrow \bigotimes_ {i \in I} F _ {i}
\]

是一个 A-代数同态，这由 (1) 得出。

对 I 的每个划分 \((\mathbf{I}_{j})_{j\in J}\) ，结合性同构

\[
\bigotimes_ {j \in J} \left(\bigotimes_ {i \in I _ {j}} \mathrm{E} _ {i}\right)\rightarrow \bigotimes_ {i \in I} \mathrm{E} _ {i}
\]

(II, § 3, 第 9 号) 也是代数同构，这由 (1) 及其定义即可得出。

当 I 是区间 \([1, n]\) （ \(\mathbf{N}\) 中的）且所有代数 \(\mathbf{E}_i\) 都等于同一个代数 \(\mathbf{E}\) 时，张量积代数 \(\bigotimes_{i \in I} \mathbf{E}_i\) 也记作 \(\mathbf{E}^{\otimes n}\) .

在本节剩余部分，我们将注意力限制在两个代数的张量积的性质上，而将把这些性质推广到任意有限族的张量积的任务留给读者。

设 E, F 为两个 A-代数；若 a（相应地 b）是 E（相应地 F）的左理想，则典范像 \(\operatorname{Im}(a \otimes b)\) （即 \(a \otimes_{A} b\) 在 \(E \otimes_{A} F\) 中的像）是 \(E \otimes_{A} F\) 的左理想；当“左理想”换成“右理想”或“双边理想”时，有类似的结果。此外：

命题 1. 设 E, F 为两个 A-代数，且 a（相应地 b）为 E（相应地 F）的一个双边理想。则典范的 A-模同构

\[
(\mathrm{E} / \mathfrak {a}) \otimes (\mathrm{E} / \mathfrak {b}) \rightarrow (\mathrm{E} \otimes \mathrm{F}) / (\operatorname{Im} (\mathfrak {a} \otimes \mathrm{F}) + \operatorname{Im} (\mathrm{E} \otimes \mathfrak {b}))
\]

（II, §3, 第 6 号, 命题 6 的推论 1）是一个代数同构。

这由 (1) 及上述引文中所给的定义得出。

推论 1. 设 E 为一个 A-代数，且 a 为 A 的一个理想。则 A-模 aE 是 E 的一个双边理想，且典范的 (A/a)-模同构

\[
(\mathbf {A} / \mathfrak {a}) \otimes_ {\mathbf {A}} \mathbf {E} \rightarrow \mathbf {E} / \mathfrak {a} \mathbf {E}
\]

是一个 (A/a)-代数同构。

推论 2. 若 \(\mathfrak{a},\mathfrak{b}\) 是 \(\mathbf{A}\) 的两个理想，则 \((\mathbf{A} / \mathfrak{a})\otimes_{\mathbf{A}}(\mathbf{A} / \mathfrak{b})\) 典范同构于 \(\mathbf{A} / (\mathfrak{a} + \mathfrak{b})\) .

推论 3. 设 E, F 为两个 A-代数，且 a 为 A 的一个包含于 F 的零化子中的理想。则 (A/a)-代数 \(E \otimes_{A} F\) 典范同构于 (E/aE) \(\otimes_{A/a} F\) .

命题 2. 设 \((\mathbf{E}_{\lambda})_{\lambda \in L}\) 和 \((\mathbf{F}_{\mu})_{\mu \in M}\) 为两个 A-代数族。典范映射（II, §3, 第 7 号）

\[
\left(\bigoplus_ {\lambda \in \mathbf {L}} \mathrm{E} _ {\lambda}\right) \otimes_ {\mathrm{A}} \left(\bigoplus_ {\mu \in \mathrm{M}} \mathrm{F} _ {\mu}\right)\rightarrow \bigoplus_ {(\lambda , \mu) \in \mathrm{L} \times \mathrm{M}} \left(\mathrm{E} _ {\lambda} \otimes_ {\mathrm{A}} \mathrm{F} _ {\mu}\right)
\]

是一个代数同构。

这直接由 II, §3, 第 7 号, 命题 7 及 \(\mathbf{E} \otimes \mathbf{F}\) 上乘法的定义得出。

命题 3. 设 A, B 为两个交换环， \(\rho: \mathbf{A} \to \mathbf{B}\) 是一个环同态，E, F 是两个 A-代数。则典范的 B-模同构

\[
\rho^ {*} (\mathrm{E}) \otimes_ {\mathbf {B}} \rho^ {*} (\mathrm{F}) \rightarrow \rho^ {*} (\mathrm{E} \otimes_ {\mathbf {A}} \mathrm{F})
\]

(II, § 5, 第 1 号, 命题 3) 是一个 B-代数同构。

命题 4. 设 A, B 为两个交换环， \(\rho: \mathbf{A} \to \mathbf{B}\) 是一个环同态，E 是一个 A-代数，F 是一个 B-代数。则典范的 A-模同构

\[
\rho_ {*} (\mathbf {F}) \otimes_ {\mathbf {A}} \mathbf {E} \rightarrow \rho_ {*} (\mathbf {F} \otimes_ {\mathbf {B}} \rho^ {*} (\mathbf {E}))
\]

(II, § 5, 第 2 号, 命题 6) 是一个 A-代数同构。

这些验证由 §1, 第 5 号 立即得出。

特别地，通过把标量环 B 限制到 A 而得到的 B \(\otimes_{\mathbf{A}} \mathbf{E}\) 上的 A-代数结构，与代数 B \(\otimes_{\mathbf{A}} \mathbf{E}\) （即 A-代数 B 与 E 的张量积）的结构相同。

最后，若 \((\mathbf{A}_i, \phi_{ji})\) 是交换环的一个正向系统， \((\mathbf{E}_i, f_{ji})\) 和 \((\mathbf{F}_i, g_{ji})\) 是两个 \(\mathbf{A}_i\) -代数的正向系统（§1, 第 6 号），且 \(\mathbf{A} = \varinjlim \mathbf{A}_i\) ，则典范的 A-模同构

\[
\varinjlim \left(\mathrm{E} _ {i} \otimes_ {\mathrm{A} _ {i}} \mathrm{F} _ {i}\right)\rightarrow (\varinjlim \mathrm{E} _ {i}) \otimes_ {\mathrm{A}} (\varinjlim \mathrm{F} _ {i})
\]

（II, § 6, 第 3 号, 命题 7）也是一个 A-代数同构，这由定义即可得出。

代数张量积的例子。（1）设 A 是一个交换环，M 和 N 是两个 A-模；典范映射

\[
\operatorname{End} _ {\mathbf {A}} (\mathbf {M}) \otimes_ {\mathbf {A}} \operatorname{End} _ {\mathbf {A}} (\mathbf {N}) \rightarrow \operatorname{End} _ {\mathbf {A}} (\mathbf {M} \otimes_ {\mathbf {A}} \mathbf {N})\tag{2}
\]

（II, § 4, 第 4 号）是一个 A-代数同态，这由 II, § 3, 第 2 号, 公式 (5) 得出。当 M 或 N 是有限生成投射 A-模时，我们知道这个同态是双射（II, § 4, 第 4 号, 命题 4）。特别地，我们重新得到两个方阵的张量积的定义。

\begin{enumerate}
\def\labelenumi{(\arabic{enumi})}
\setcounter{enumi}{1}
\tightlist
\item
  设 \(S, T\) 是两个幺半群， \(A^{(S)}\) 和 \(A^{(T)}\) 是幺半群 \(S\) 和 \(T\) 在环 \(A\) 上的代数（III, § 2, 第 6 号）；则存在一个典范的 \(A\) -代数同构
\end{enumerate}

\[
\mathbf {A} ^ {(\mathbf {S})} \otimes_ {\mathbf {A}} \mathbf {A} ^ {(\mathbf {T})} \rightarrow \mathbf {A} ^ {(\mathbf {S} \times \mathbf {T})}.\tag{3}
\]

元素 \(e_s \otimes e_t\) （相应地， \(e_{(s,t)}\) ），其中 \(s\) 遍历 \(S\) ， \(t\) 遍历 \(T\) ，由 II, §3, 第 7 号, 命题 7 的推论 2，构成 \(A^{(S)} \otimes_A A^{(T)}\) 的一个基（相应地，构成 \(A^{(S \times T)}\) 的一个基）；把 \(e_s \otimes e_t\) 映到 \(e_{(s,t)}\) 就得到所求的同构，且由定义可知它确是代数同构。

